\section{Martingales}

	Remember that for a random variable $X$ and $\sigma$-algebras $\mathscr{F}_1, \mathscr{F}_2$ s.t. $\mathscr{F}_1 \subset \mathscr{F}_2$, by the tower property, 
	\begin{equation}
		\mathbb{E}[ \mathbb{E}[X \mid \mathscr{F}_2] \mid \mathscr{F}_1] = \mathbb{E}[X \mid \mathscr{F}_1] 
	\end{equation}

	If we exchange the two, the smaller one wins. but for arbitrary $\sigma$-algebras, equality does not hold.\footnote{No direct connection to Fubini's theorem.}
	\begin{equation}
		\mathbb{E}[ \mathbb{E}[X \mid \mathscr{F}_2] \mid \mathscr{F}_1] \neq \mathbb{E}[ \mathbb{E}[X \mid \mathscr{F}_1] \mid \mathscr{F}_2]
	\end{equation}

	It is useful to consider when $\mathscr{F}_2 \not\subset \mathscr{F}_1$ and $\mathscr{F}_1 \not\subset \mathscr{F}_2$. 

	\begin{example}
		The prototypical example is $\Omega = \{a, b, c\}$ with $\mathbb{P}(\{a\}) = \mathbb{P}(\{b\}) = \mathbb{P}(\{c\}) = \frac{1}{3}$. Then, $X(a) = 1, X(b) = 2, X(c) = 3$. Consider the following $\sigma$-algebras: 
		\begin{enumerate}
			\item $\mathscr{F}_1$ generated by $\{a, b\}, \{c\}$. 
			\item $\mathscr{F}_2$ generated by $\{a\}, \{b, c\}$. 
		\end{enumerate}
		Then, 
		\begin{align}
			\mathbb{E}[X \mid \mathscr{F}_1](\omega) & = \begin{cases} 
				1.5 & \text{ if } \omega \in \{a, b\} \\ 
				3 & \text{ if } \omega \in \{c\} 
			\end{cases} \implies \mathbb{E}[ \mathbb{E}[X \mid \mathscr{F}_1] \mid \mathscr{F}_2] = 
			\begin{cases} 
				1.5 & \text{ if } \omega \in \{a\} \\ 
				2.75 & \text{ if } \omega \in \{b, c\} 
			\end{cases} \\ 
			\mathbb{E}[X \mid \mathscr{F}_2](\omega) & = \begin{cases} 
				1 & \text{ if } \omega \in \{a\} \\ 
				2.5 & \text{ if } \omega \in \{b, c\} 
			\end{cases} \implies \mathbb{E}[ \mathbb{E}[X \mid \mathscr{F}_2] \mid \mathscr{F}_1] = 
			\begin{cases} 
				1.75 & \text{ if } \omega \in \{a, b\} \\ 
				2.5 & \text{ if } \omega \in \{c\} 
			\end{cases}
		\end{align}
	\end{example}

	\begin{definition}[Filtration]
		A collection of increasing $\sigma$-algebras $\mathscr{F} = (\mathscr{F}_n)_{n=0}^\infty$ with $\mathscr{F}_n \subset \mathscr{F}_{n+1}$ is a \textbf{filtration}. A sequence of random variables $(X_n)_{n=0}^\infty$ is \textbf{adapted} to $\mathscr{F}$ if for every $n$, $X_n$ is $\mathscr{F}_n$-measurable. 
	\end{definition}

	\begin{definition}[Martingale]
	  If $X_n$ further satisfies 
		\begin{enumerate}
			\item $\mathbb{E}[|X_n|] < +\infty$ (so that we don't run into any problems in cord exp)
			\item $\mathbb{E}[X_n \mid \mathscr{F}_{n-1}] = X_{n-1}$. 
		\end{enumerate}
		Then $(X_n)$ is called a \textbf{martingale}. If instead of (2), we have 
		\begin{enumerate}
		  \item $\mathbb{E}[X_n \mid \mathscr{F}_{n-1}] \leq X_{n-1}$, then $(X_n)$ is a \textbf{supermartingale}. 
		  \item $\mathbb{E}[X_n \mid \mathscr{F}_{n-1}] \geq X_{n-1}$, then $(X_n)$ is a \textbf{submartingale}. 
		\end{enumerate}
	\end{definition}

	\begin{example}[Random Walk]
	  Let $Y_1, \ldots$ be iid random variables and define 
		\begin{equation}
			S_n = S_0 + \sum_{i=1}^n Y_i
		\end{equation}
		where $S_0$ is a constant. To specify a filtration, a canonical way of constructing it is to just let it be generated by the random variables. So, 
		\begin{equation}
			\mathscr{F}_0 = \{\empty,  \Omega\}, \qquad \mathscr{F}_n = \sigma(Y_1, \ldots, Y_n)
		\end{equation}
		For (1) to hold, just assume $\mathbb{E}[|Y_i|] < +\infty \implies \mathbb{E}[|S_n|] < +\infty$ by triangle inequality. So, by our construction, we get 
		\begin{equation}
			\mathbb{E}[S_n \mid \mathscr{F}_{n-1}] = \underbrace{S_0 + \sum_{i=1}^{n-1} Y_i}_{= S_{n-1}} + \underbrace{\mathbb{E}[Y_n \mid \mathscr{F}_{n-1}]}_{= \mathbb{E}[Y_n]} 
		\end{equation}
		where $\mathbb{E}[Y_n \mid \mathscr{F}_{n-1}] = \mathbb{E}[Y_n]$ since $Y_n$ is independent of $\mathscr{F}_{n-1}$. So, 
		\begin{enumerate}
			\item if $\mathbb{E}[Y_n] = 0$, $S_n$ is a martingale. 
			\item if $\mathbb{E}[Y_n] > 0$, $S_n$ is a submartingale. 
			\item if $\mathbb{E}[Y_n] < 0$, $S_n$ is a supermartingale. 
		\end{enumerate}
	\end{example}

	\begin{example}
	  Now let 
		\begin{equation}
			M_n = \prod_{i=1}^n Y_i, \qquad Y_i \geq 0
		\end{equation}
		and assume $\mathbb{E}[Y_i]^2 < +\infty$. Then, by Cauchy-Schwarz, $\mathbb{E}[|M_n|] < +\infty$. Then, 
		\begin{equation}
			\mathbb{E}[M_n \mid \mathscr{F}_{n-1}] = \bigg( \prod_{i=1}^n Y_i \bigg) \mathbb{E}[Y_n] = M_{n-1} \mathbb{E}[Y_n]
		\end{equation}
		So, 
		\begin{enumerate}
			\item if $\mathbb{E}[Y_n] = 1$, $S_n$ is a martingale. 
			\item if $\mathbb{E}[Y_n] > 1$, $S_n$ is a submartingale. 
			\item if $\mathbb{E}[Y_n] < 1$, $S_n$ is a supermartingale. 
		\end{enumerate}
	\end{example}

	Now for some abstract properties of martingales. 

	\begin{lemma} 
		If $X_n$ is a supermartingale, then for all $n > m$, $\mathbb{E}[X_n \mid \mathscr{F}_m] \leq X_m$.
	\end{lemma}
	\begin{proof}
	  Induct on $n$ and tower property. 
		\begin{align}
			\mathbb{E}[X_n \mid \mathscr{F}_m] & = \mathbb{E}[\mathbb{E}[X_n \mid \mathscr{F}_{n-1}] \mid \mathscr{F}_m ] && \tag{Tower Property} \\ 
																				 & \leq \mathbb{E}[X_{n-1} \mid \mathscr{F}_m ] && \tag{By Definition} \\ 
																				 & \leq X_m && \tag{Inductive Hypothesis}
		\end{align}
	\end{proof}

	\begin{corollary}
	  List. 
		\begin{enumerate}
			\item If $X_n$ is a submartingale, then for all $n > m$, $\mathbb{E}[X_n \mid \mathscr{F}_m] \geq X_m$. 
			\item If $X_n$ is a  martingale, then for all $n > m$, $\mathbb{E}[X_n \mid \mathscr{F}_m] = X_m$. 
		\end{enumerate}
	\end{corollary}
	\begin{proof}
		We use the common ``symmetry trick.''
		\begin{enumerate}
		  \item $-X_n$ is a supermartingale and use the lemma. 
			\item $X_n$ is both a supermartingale and submartingale. 
		\end{enumerate}
	\end{proof}

	\begin{definition}[Predictable Sequence]
		For a filtration $(\mathscr{F}_n)_{n=0}^\infty$, a sequence of random variables $(H_n)_{n=1}^\infty$ is a \textbf{predictable sequence} if $H_n$ is measurable with respect to $\mathscr{F}_{n-1}$. 
	\end{definition}

	It is immediate that predictable sequences are adapted, since the $\sigma$-algebras are increasing in a filtration, i.e. $H_n$ is $\mathscr{F}_{n-1}$-measurable which implies $H_n$ is $\mathscr{F}_n$ measurable. Adapated sequences are not necessarily predictable. 

	\begin{example}[Gambling]
	  At each round $r$, you bet \$x. With probability $p$, you win an additional \$x and with probability $1-p$, you lose \$x. Let 
		\begin{enumerate}
		  \item $X_n = $ not winning after $n$ rounds, if you bet \$1 each time. 
			\item $H_n = $ amount of money you bet in round $n+1$. So, makes sense to be $\mathscr{F}_{n+1}$-measurable. So, this is an example of a predictable sequence. 
			\item Now, our net earnings after $n$ rounds of the game is 
				\begin{equation}
					(H \cdot X)_n = \sum_{m=1}^n H_m (X_m - X_{m-1})
				\end{equation}
				where $X_m - X_{m-1} \in \{\pm 1\}$. 
		\end{enumerate}
		Now the strategy might be to double down if you lose, and keep betting \$1 if you win, so 
		\begin{equation}
			H_n = \begin{cases} 
				2 H_{n-1} & \text{ if } X_{n-1} - X_{n-2} = -1 \\ 
				1 & \text{ if } X_{n-1} - X_{n-2} = \pm 1
			\end{cases} 
		\end{equation}
		This betting strategy is called a martingale. Remember that if we had infinite money, we will always win. More generally, if $(H \cdot X)_n$ is a submartingale, then $\mathbb{E}[(H \cdot X)_n] \geq 0$. So, we should consider cases which leads to $(H \cdot X)_n$ being a submartingale. 
	\end{example}

	\begin{theorem}
		Let $X_n$, $n \geq 0$ be a supermartingale. If $H_n \geq 0$ is predictable and each $H_n$ is bounded (pointwise), then $(H \cdot X)_n$ is a super-martingale. 
	\end{theorem}
	\begin{proof}
	  We have 
		\begin{align}
			\mathbb{E}[(H \cdot X)_{n+1} \mid \mathscr{F}_n] & = (H \cdot X)_n + \mathbb{E}[H_{n+1} (X_{n+1} - X_n) \mid \mathscr{F}_n ] \\ 
																											 & = (H \cdot X)_n + H_{n+1} \mathbb{E}[\underbrace{X_{n+1} - X_n}_{\leq 0} \mid \mathscr{F}_n ] \\ &
																											& \leq (H \cdot X)_n 
		\end{align}
	\end{proof}

	\begin{definition}[Stopping Time]
	  We call a natural-valued random variable $N$ a \textbf{stopping time} if 
		\begin{equation}
			\{N = n\} = N^{-1} (n) \in \mathscr{F}_n, \qquad ( \mathscr{F}_k \subseteq \mathscr{F}_n \; \forall k \leq n \implies \{N = k\} \in \mathscr{F}_n )
		\end{equation}
		This is abstract, but just think of it as a strategy of stopping a betting game. 
	\end{definition}

	\begin{theorem}
		If $N$ is a stopping time, and $X_n$ is a supermartingale, then $X_{N \wedge n}$ (where $N \wedge n$) is a supermartingale. This is a stopped version. 
	\end{theorem}
	\begin{proof}
	  We can write 
		\begin{equation}
			X_{N \cap n} = \bigg( \sum_{m=1}^n H_m (X_m - X_{m-1}) \bigg) + X_0
		\end{equation}
		where $H_m = \mathbbm{1}(N \geq m)$. Check that it is predictable, and then conclude by theorem 1. 
	\end{proof}

	\begin{example}
		Suppose $(X_n)_{n \geq 0}$ is a submartingale. Let $a < b$, $N_0 = -1$. For $k \geq 1$, let 
		\begin{enumerate}
			\item $N_{2k-1} = \inf\{m > N_{2k-2} \colon X_m \leq a \}$, i.e. when are you going below $a$? 
			\item $N_{2k} = \inf \{m > N_{2k-1} \colon X_m \geq b \}$ , i.e. when are you going above $b$? 
		\end{enumerate}
		We claim that $N_j$ are stopping times. We call $U_n = \sup \{k \colon N_{2k} \leq n\}$ as the number of upcrossings completed by time $n$. Next, we have the upcrossing inequality, which is used the prove the martingale convergence theorem.  
	\end{example}

	\begin{theorem}[Upcrossing Inequality]
		If $(X_m)_{m \geq 0}$ is a submartingale, then 
		\begin{equation}
			(b - a) \mathbb{E}[U_n] \leq \mathbb{E}[(X_n - a)^+] - \mathbb{E}[(X_0 - a)^+]
		\end{equation}
	\end{theorem}

	The martingale convergence theorem is like an analogous version of the Bolzano-Weierstrass theorem---that a monotone sequence bounded above converges. 

